% =====================================================
% PREÁMBULO
%=======================================================

% DEFINIR EL TIPO DE DOCUMENTO {}
% DEFINIR EL TAMAÑO DE LETRA []
% DEFINIR LAS DIMENSIONES []: letterpaper | a4paper
\documentclass[a4paper,12pt]{article}
\usepackage{graphicx}

% Lenguaje
\usepackage[spanish]{babel}
\usepackage[utf8]{inputenc}

% Margenes
\usepackage{geometry}
\geometry{
    top=3cm,
    bottom=3cm,
    left=2.5cm,
    right=2.5cm
}

% Texto Justificado
\usepackage{ragged2e}
\justifying

% Interlineado
\usepackage{setspace}
\setstretch{1.5}

% Sangría
\setlength{\parindent}{1.5cm}

% Metadatos del documento
\title{Taller de Latex - Sesión 1}
\author{Facundo Cabral}
\date{Septiembre de 2026}

% Paquetes Adicionales
\usepackage{amsfonts}
\usepackage{amsmath}
\usepackage{amssymb}
\usepackage{amsthm}

\usepackage{xcolor}

\usepackage{enumitem}

% Funciones Creadas
\newcommand{\azul}[1]{\textcolor{blue}{#1}}

\newcommand{\der}[2]{\frac{\partial #1}{\partial #2}}

% =====================================================
% DOCUMENTO
%=======================================================

\begin{document}

\maketitle

\newpage

\tableofcontents

\newpage

%=======================================================
\section{Tipos de Letras}
%=======================================================

Este es un texto normal.

\textbf{Este es un texto en negrita.}

\textit{Este es un texto en itálica.}

\textsl{Este texto está inclinado.}

\textsc{Este texto está en versalitas.}

\textcolor{blue}{Este texto está en azul.}

\azul{Este texto está en azul con función creada}

%=======================================================
\section{Listas y Enumeraciones}
%=======================================================

% Sin sangría este párrafo
\noindent A continuación, se muestra la lista de pendientes:

\begin{itemize}[leftmargin=*, nosep]
    \item Pendiente 1
    \item[$\bullet$] Pendiente 2
    \item[$\circ$] Pendiente 3
\end{itemize}

\vspace*{0.5cm}

A continuación, lista enumerada:

\begin{enumerate}
    \item Pendiente 1
    \item Pendiente 2
\end{enumerate}

\newpage

\section{Matemática}

Un consumidor racional elige la cesta de bienes que maximiza su utilidad sujeto a la restriccion presupuestaria. Consideramos un ingreso $I>0$ que se enfrenta a un vector de precios $(p_x,p_y) \in \mathbb{R}^2_{++}$ para los bienes $x$ e $y$.

\vspace{0.5cm}

La propiedad de transitividad de las preferencias $$u(x) \geq u(y) \iff x \succsim y$$

$$u(x,y)=A\cdot x^\alpha \cdot y^{1-\alpha} \hspace{1cm} \forall A>0, \alpha \in (0,1)$$

\vspace{-1cm}

\begin{align*}
    \max_{x,y} \hspace{0.25cm} & u(x,y) \\
    & \text{s.a. } \hspace{0.25cm} p_xx+p_yy = I
\end{align*}

\vspace{-1cm}

\begin{align*}
    a & = \beta\cdot 5 \\
    a/\beta & = 5
\end{align*}


$$u(x,y)= x^\alpha \cdot y^{1-\alpha}$$

\vspace{-1cm}

\begin{align*}
    \mathcal{L}(x,y,\lambda) & = x^\alpha \cdot y^{1-\alpha} + \lambda(I-p_xx-p_yy)\\
    \der{\mathcal{L}}{x} & = \alpha x^{\alpha-1}y^{1-\alpha} - \lambda p_x =0\\
    \der{\mathcal{L}}{y} & = (1-\alpha) x^{\alpha}y^{-\alpha} - \lambda p_y=0 \\
    \der{\mathcal{L}}{\lambda } & = I-p_xx-p_yy=0 \\
\end{align*}

\begin{equation}
    \boxed{
        x^*(p_x,p_y,I)=\frac{\alpha I}{p_x} \hspace{1.5cm}y^*(p_x,p_y,I)=\frac{(1-\alpha)I}{p_y}
    }
\end{equation}

\[
    \begin{bmatrix}
        0   & p_x   & p_y \\
        p_x & U_{xx}    & U_{xy} \\
        p_y & U_{yx}    & U_{yy}
    \end{bmatrix}_{3\times3}
    \begin{bmatrix}
        \partial \lambda \\
         \partial x \\
         \partial y
    \end{bmatrix}_{3\times 1} =
    \begin{bmatrix}
        0 \\
        0 \\
        0
    \end{bmatrix}
\]

$$\sum_{t=0}^{T}\beta^t \cdot U(c_t)$$

$$\int_{0}^{\infty}e^{-rt}\pi_tdt$$

$$\lim_{n \to \infty} \left\{1+\frac{1}{n}\right\}^n=e$$


\end{document}
